% Tento soubor nahraďte vlastním souborem s přílohami (nadpisy níže jsou pouze pro příklad)

% Pro kompilaci po částech (viz projekt.tex), nutno odkomentovat a upravit
%\documentclass[../projekt.tex]{subfiles}
%\begin{document}

% Umístění obsahu paměťového média do příloh je vhodné konzultovat s vedoucím
\chapter{Obsah paměťového média}
Paměťové médium obsahuje následující soubory a složky, které jsou uložené v~adresáři s~názvem \uv{xoravo01-Identifikace-zvere-pomoci-AI}:

\begin{itemize}
	\item \textbf{README.md} -- Základní informace o~projektu, jeho struktuře, spuštění a~ovládání programu.
	
	\item \textbf{thesis/} -- Zdrojové soubory textu v~\LaTeX{}u.
	
	\item \textbf{final\_thesis.pdf} -- Text bakalářské práce ve formátu PDF.
	
	\item \textbf{implementation/} -- Zdrojové kódy, modely a~data.
	\begin{itemize}
		\item \textbf{main.py} -- Spouštěcí skript.
		\item \textbf{detection/} -- Zdrojové kódy související s~detekcí.
		\item \textbf{identification/} -- Zdrojové kódy související s~identifikací.
		\item \textbf{data/} -- Data použitá pro trénování a~testování modelů.
		\item \textbf{models/} -- Natrénované modely neuronových sítí.
		\item \textbf{db/} -- Zdrojové kódy pro práci s~databází a~záložní (\textit{backup}) soubory.
		\item \textbf{app/} -- Uživatelské rozhraní aplikace (GUI).
		\item \textbf{documentation.pdf} -- Technická dokumentace popisující strukturu a~funkce programu.
		\item \textbf{requirements.txt} -- Seznam knihoven a~závislostí potřebných pro spuštění.
	\end{itemize}
	
\end{itemize}


% Pro kompilaci po částech (viz projekt.tex) nutno odkomentovat
%\end{document}
